% chapterintro
% ==============================
% 第11章 导言：11.0-chinese.md
本章的主要内容有：

\begin{itemize}
\item 生成源代码和二进制包

\item 使用CMake/pybind11构建的C++/Python项目，通过PyPI发布

\item 使用CMake/CFFI构建C/Fortran/Python项目，通过PyPI发布

\item 以Conda包的形式发布一个简单的项目

\item 将Conda包作为依赖项发布给项目

\end{itemize}

目前为止，已经从源代码编译并安装了软件包——这意味着可以通过Git获取项目，并手动执行配置、构建、测试和安装。然而，在实际中，软件包通常是使用管理器来安装的，比如Apt、DNF、Pacman、pip和Conda。我们需要以各种格式发布我们的代码项目——作为源文件或二进制安装程序。

下图显示了使用CMake的项目的各个阶段，我们现在方案就是其中的打包时(PI)：

\myGraphic{0.9}{images/preface/2.png}{}

本章中，我们将探讨不同的打包策略。首先，讨论使用CMake中的工具CPack进行打包，还提供打包和上传CMake项目到Python包索引(PyPI, \texthref{https://pypi.org})和Anaconda云(https://anaconda.org )的方法，这些都是通过包管理器pip和Conda (https://conda.io/docs/ )分发包的平台。对于PyPI，我们将演示如何打包和分发混合C++/Python或C/Fortran/Python的项目。对于Conda，我们将展示如何对依赖于其他库的C++项目进行打包。


% --- recipe 1: 11.1-chinese.md ---
\section{11.1 生成源代码和二进制包}

\begin{myInfoBox}
\textbf{NOTE}：此示例代码可以在 \texthref{https://github.com/dev-cafe/cmake-cookbook/tree/v1.0/chapter-11/recipe-01} 中找到。该示例在CMake 3.6版(或更高版本)中是有效的，并且已经在GNU/Linux、macOS和Windows上进行过测试。
\end{myInfoBox}

如果代码是开源的，用户将能够下载项目的源代码，并使用完全定制的CMake脚本自行构建。当然，打包操作也可以使用脚本完成，但是CPack提供了更简单和可移植的替代方案。本示例将指导您创建一些包:

\begin{itemize}
\item \textbf{源代码包}：可以将源代码直接压缩成归档文件，进行发送。用户将不必为特定的版本控制系统操心。

\item \textbf{二进制包}：工具将新构建的目标以打包的方式到归档文件中。这个功能非常有用，但可能不够健壮，无法发布库和可执行程序。

\item \textbf{平台原生的二进制安装}：CPack能够以许多不同的格式生成二进制安装程序，因此可以将软件发布到不同的平台。我们将展示如何生成安装程序:

\begin{itemize}
\item 基于Debian的GNU/Linux发行版的\texttt{.deb}格式： \texthref{https://manpages.debian.org/unstable/dpkg-dev/deb.5.en.html}

\item 基于Red Hat的GNU/Linux发行版的\texttt{.rpm}格式： \texthref{http://rpm.org/}

\item macOS包的\texttt{.dmg}格式: \texthref{https://developer.apple.com/library/archive/documentation/CoreFoundation/Conceptual/CFBundles/BundleTypes/BundleTypes.html}

\item Windows的NSIS格式: \texthref{http://nsis.sourceforge.net/Main\_Page}

\end{itemize}

\end{itemize}

\mySmallsection{准备工作}

我们将使用第10章第3节的示例，项目树由以下目录和文件组成:

\begin{codebox}[shell]
.
├── cmake
│    ├── coffee.icns
│    ├── Info.plist.in
│    └── messageConfig.cmake.in
├── CMakeCPack.cmake
├── CMakeLists.txt
├── INSTALL.md
├── LICENSE
├── src
│    ├── CMakeLists.txt
│    ├── hello-world.cpp
│    ├── Message.cpp
│    └── Message.hpp
└── tests
    ├── CMakeLists.txt
    └── use_target
        ├── CMakeLists.txt
        └── use_message.cpp
\end{codebox}

由于本示例的重点是使用CPack，所以不会讨论源码。我们只会在\texttt{CMakeCPack.cmake}中添加打包指令。此外，还添加了\texttt{INSTALL.md}和\texttt{LICENSE}文件：打包要求需要包含安装说明和项目许可信息。

\mySmallsection{具体实施}

让我们看看需要添加到这个项目中的打包指令。我们将在\texttt{ CMakeCPack.cmake }中收集它们，并在在\texttt{CMakeLists.txt}的末尾包含这个模块\texttt{include(cmakecpackage.cmake)}:

\begin{enumerate}
\item 我们声明包的名称，与项目的名称相同，因此我们使用\texttt{PROJECT\_NAME}的CMake变量:

\begin{codebox}[cmake]
set(CPACK_PACKAGE_NAME "${PROJECT_NAME}")
\end{codebox}

\item 声明包的供应商：

\begin{codebox}[cmake]
set(CPACK_PACKAGE_VENDOR "CMake Cookbook")
\end{codebox}

\item 打包的源代码将包括一个描述文件。这是带有安装说明的纯文本文件:

\begin{codebox}[cmake]
set(CPACK_PACKAGE_DESCRIPTION_FILE "${PROJECT_SOURCE_DIR}/INSTALL.md")
\end{codebox}

\item 还添加了一个包的描述:

\begin{codebox}[cmake]
set(CPACK_PACKAGE_DESCRIPTION_SUMMARY "message: a small messaging library")
\end{codebox}

\item 许可证文件也将包括在包中:

\begin{codebox}[cmake]
set(CPACK_RESOURCE_FILE_LICENSE "${PROJECT_SOURCE_DIR}/LICENSE")
\end{codebox}

\item 从发布包中安装时，文件将放在\texttt{/opt/recipe-01}目录下:

\begin{codebox}[cmake]
set(CPACK_PACKAGING_INSTALL_PREFIX "/opt/${PROJECT_NAME}")
\end{codebox}

\item CPack所需的主要、次要和补丁版本:

\begin{codebox}[cmake]
set(CPACK_PACKAGE_VERSION_MAJOR "${PROJECT_VERSION_MAJOR}")
set(CPACK_PACKAGE_VERSION_MINOR "${PROJECT_VERSION_MINOR}")
set(CPACK_PACKAGE_VERSION_PATCH "${PROJECT_VERSION_PATCH}")
\end{codebox}

\item 设置了在包装的时候需要忽略的文件列表和目录:

\begin{codebox}[cmake]
set(CPACK_SOURCE_IGNORE_FILES "${PROJECT_BINARY_DIR};/.git/;.gitignore")
\end{codebox}

\item 列出了源代码归档的打包生成器——在我们的例子中是\texttt{ZIP}，用于生成\texttt{.ZIP}归档，\texttt{TGZ}用于\texttt{.tar.gz}归档:

\begin{codebox}[cmake]
set(CPACK_SOURCE_GENERATOR "ZIP;TGZ")
\end{codebox}

\item 我们还列出了二进制存档生成器:

\begin{codebox}[cmake]
set(CPACK_GENERATOR "ZIP;TGZ")
\end{codebox}

\item 现在也可声明平台原生二进制安装程序，从DEB和RPM包生成器开始，不过只适用于GNU/Linux:

\begin{codebox}[cmake]
if(UNIX)
  if(CMAKE_SYSTEM_NAME MATCHES Linux)
    list(APPEND CPACK_GENERATOR "DEB")
    set(CPACK_DEBIAN_PACKAGE_MAINTAINER "robertodr")
    set(CPACK_DEBIAN_PACKAGE_SECTION "devel")
    set(CPACK_DEBIAN_PACKAGE_DEPENDS "uuid-dev")
    list(APPEND CPACK_GENERATOR "RPM")
    set(CPACK_RPM_PACKAGE_RELEASE "1")
    set(CPACK_RPM_PACKAGE_LICENSE "MIT")
    set(CPACK_RPM_PACKAGE_REQUIRES "uuid-devel")
  endif()
endif()
\end{codebox}

\item 如果我们在Windows上，我们会想要生成一个NSIS安装程序:

\begin{codebox}[cmake]
if(WIN32 OR MINGW)
  list(APPEND CPACK_GENERATOR "NSIS")
  set(CPACK_NSIS_PACKAGE_NAME "message")
  set(CPACK_NSIS_CONTACT "robertdr")
  set(CPACK_NSIS_ENABLE_UNINSTALL_BEFORE_INSTALL ON)
endif()
\end{codebox}

\item 另一方面，在macOS上，bundle包是我们的安装程序的选择:

\begin{codebox}[cmake]
if(APPLE)
  list(APPEND CPACK_GENERATOR "Bundle")
  set(CPACK_BUNDLE_NAME "message")
  configure_file(${PROJECT_SOURCE_DIR}/cmake/Info.plist.in Info.plist @ONLY)
  set(CPACK_BUNDLE_PLIST ${CMAKE_CURRENT_BINARY_DIR}/Info.plist)
  set(CPACK_BUNDLE_ICON ${PROJECT_SOURCE_DIR}/cmake/coffee.icns)
endif()
\end{codebox}

\item 我们在现有系统的包装生成器上，向用户打印一条信息:

\begin{codebox}[cmake]
message(STATUS "CPack generators: ${CPACK_GENERATOR}")
\end{codebox}

\item 最后，我们包括了\texttt{CPack.cmake}标准模块。这将向构建系统添加一个包和一个\texttt{package\_source}目标:

\begin{codebox}[cmake]
include(CPack)
\end{codebox}

\end{enumerate}

现在来配置这个项目：

\begin{codebox}[shell]
$ mkdir -p build
$ cd build
$ cmake .. 
\end{codebox}

使用下面的命令，我们可以列出可用的目标(示例输出是在使用Unix Makefile作为生成器的GNU/Linux系统上获得的):

\begin{codebox}[shell]
$ cmake --build . --target help

The following are some of the valid targets for this Makefile:
... all (the default if no target is provided)
... clean
... depend
... install/strip
... install
... package_source
... package
... install/local
... test
... list_install_components
... edit_cache
... rebuild_cache
... hello- world
... message
\end{codebox}

我们可以看到\texttt{package}和\texttt{package\_source}目标是可用的。可以使用以下命令生成源包:

\begin{codebox}[shell]
$ cmake --build . --target package_source

Run CPack packaging tool for source...
CPack: Create package using ZIP
CPack: Install projects
CPack: - Install directory: /home/user/cmake-cookbook/chapter-11/recipe-01/cxx-example
CPack: Create package
CPack: - package: /home/user/cmake-cookbook/chapter- 11/recipe-01/cxx-example/build/recipe-01-1.0.0-Source.zip generated.
CPack: Create package using TGZ
CPack: Install projects
CPack: - Install directory: /home/user/cmake-cookbook/chapter- 11/recipe-01/cxx-example
CPack: Create package
CPack: - package: /home/user/cmake-cookbook/chapter-11/recipe-01/cxx-example/build/recipe-01- 1.0.0-Source.tar.gz generated. 
\end{codebox}

同样，也可以构建二进制包:

\begin{codebox}[shell]
$ cmake --build . --target package message-1.0.0-Linux.deb
\end{codebox}

例子中，最后得到了以下二进制包:

\begin{codebox}[shell]
message-1.0.0-Linux.rpm
message-1.0.0-Linux.tar.gz
message-1.0.0-Linux.zip
\end{codebox}

\mySmallsection{工作原理}

CPack可用于生成用于分发的包。生成构建系统时，我们在\texttt{CMakeCPack.cmake}中列出了CPack指令，用于在构建目录下生成\texttt{ CPackConfig.cmake}。当运行以\texttt{package}或\texttt{package\_source}目标的CMake命令时，CPack会自动调用，参数是自动生成的配置文件。实际上，这两个新目标是对CPack简单规则的使用。与CMake一样，CPack也有生成器的概念。CMake上下文中的生成器是用于生成本地构建脚本的工具，例如Unix Makefile或Visual Studio项目文件，而CPack上下文中的生成器是用于打包的工具。我们列出了这些变量，并对不同的平台进行了特别的关注，为源包和二进制包定义了\texttt{CPACK\_SOURCE\_GENERATOR}和\texttt{CPACK\_GENERATOR}变量。因此，\texttt{DEB}包生成器将调用\texttt{Debian}打包实用程序，而\texttt{TGZ}生成器将调用给定平台上的归档工具。我们可以直接在\texttt{build}目录中调用CPack，并选择要与\texttt{ CPackConfig.cmake}0命令行选项一起使用的生成器。\texttt{ CPackConfig.cmake}1包可以通过以下步骤生成:

\begin{codebox}[shell]
$ cd build
$ cpack -G RPM

CPack: Create package using RPM
CPack: Install projects
CPack: - Run preinstall target for: recipe-01
CPack: - Install project: recipe-01
CPack: Create package
CPackRPM: Will use GENERATED spec file: /home/user/cmake-cookbook/chapter-11/recipe-01/cxx-example/build/_CPack_Packages/Linux/RPM/SPECS/recipe-01.spec
CPack: - package: /home/user/cmake-cookbook/chapter-11/recipe-01/cxx-example/build/recipe-01-1.0.0-Linux.rpm generated.
\end{codebox}

对于任何发行版，无论是源代码还是二进制文件，我们只需要打包用户需要的内容，因此整个构建目录和其他与版本控制相关的文件，都必须从要打包的文件列表中排除。我们的例子中，排除列表使用下面的命令声明：

\begin{codebox}[cmake]
set(CPACK_SOURCE_IGNORE_FILES "${PROJECT_BINARY_DIR};/.git/;.gitignore")
\end{codebox}

我们还需要指定包的基本信息，例如：名称、简短描述和版本。这个信息是通过CMake变量设置的，当包含相应的模块时，CMake变量被传递给CPack。

\begin{myInfoBox}
\textbf{NOTE}：由于CMake 3.9中的\texttt{project()}命令接受\texttt{DESCRIPTION}字段，该字段带有一个描述项目的短字符串。CMake将设置一个\texttt{PROJECT\_DESCRIPTION}，可以用它来重置\texttt{CPACK\_PACKAGE\_DESCRIPTION\_SUMMARY}。
\end{myInfoBox}

让我们详细看看，可以为示例项目生成的不同类型包的说明。

\mySmallsection{打包源码}

我们的示例中，决定对源存档使用\texttt{TGZ}和\texttt{ZIP}生成器。这些文件将分别生成\texttt{.tar.gz}和\texttt{.zip}压缩文件。我们可以检查生成的\texttt{.tar.gz}文件的内容:

\begin{codebox}[shell]
$ tar tzf recipe-01-1.0.0-Source.tar.gz

recipe-01-1.0.0-Source/opt/
recipe-01-1.0.0-Source/opt/recipe-01/
recipe-01-1.0.0-Source/opt/recipe-01/cmake/
recipe-01-1.0.0-Source/opt/recipe-01/cmake/coffee.icns
recipe-01-1.0.0-Source/opt/recipe-01/cmake/Info.plist.in
recipe-01-1.0.0-Source/opt/recipe-01/cmake/messageConfig.cmake.in
recipe-01-1.0.0-Source/opt/recipe-01/CMakeLists.txt
recipe-01-1.0.0-Source/opt/recipe-01/src/
recipe-01-1.0.0-Source/opt/recipe-01/src/Message.hpp
recipe-01-1.0.0-Source/opt/recipe-01/src/CMakeLists.txt
recipe-01-1.0.0-Source/opt/recipe-01/src/Message.cpp
recipe-01-1.0.0-Source/opt/recipe-01/src/hello-world.cpp
recipe-01-1.0.0-Source/opt/recipe-01/LICENSE
recipe-01-1.0.0-Source/opt/recipe-01/tests/
recipe-01-1.0.0-Source/opt/recipe-01/tests/CMakeLists.txt
recipe-01-1.0.0-Source/opt/recipe-01/tests/use_target/
recipe-01-1.0.0-Source/opt/recipe-01/tests/use_target/CMakeLists.txt
recipe-01-1.0.0-Source/opt/recipe-01/tests/use_target/use_message.cpp
recipe-01-1.0.0-Source/opt/recipe-01/INSTALL.md
\end{codebox}

与预期相同，只包含源码树的内容。注意\texttt{INSTALL.md }和\texttt{LICENSE}文件也包括在内，可以通过\texttt{CPACK\_PACKAGE\_DESCRIPTION\_FILE}和\texttt{CPACK\_RESOURCE\_FILE\_LICENSE}变量指定。

\begin{myInfoBox}
\textbf{NOTE}：Visual Studio生成器无法解析\texttt{package\_source}目标:\texthref{https://gitlab.kitware.com/cmake/cmake/issues/13058。}
\end{myInfoBox}

\mySmallsection{二进制包}

创建二进制存档时，CPack将打包\texttt{CMakeCPack.cmake}中描述的目标的内容。因此，在我们的示例中，hello-world可执行文件、消息动态库以及相应的头文件都将以\texttt{.tar.gz}和\texttt{.zip}的格式打包。此外，还将打包CMake配置文件。这对于需要链接到我们的库的其他项目非常有用。包中使用的安装目录可能与从构建树中安装项目时使用的前缀不同，可以使用\texttt{CPACK\_PACKAGING\_INSTALL\_PREFIX}变量来实现这一点。我们的示例中，我们将它设置为系统上的特定位置:\texttt{/opt/recipe-01}。

\begin{codebox}[shell]
$ tar tzf recipe-01-1.0.0-Linux.tar.gz

recipe-01- 1.0.0-Linux/opt/
recipe-01-1.0.0-Linux/opt/recipe-01/
recipe-01-1.0.0- Linux/opt/recipe-01/bin/
recipe-01-1.0.0-Linux/opt/recipe-01/bin/hello- world
recipe-01-1.0.0-Linux/opt/recipe-01/share/
recipe-01-1.0.0- Linux/opt/recipe-01/share/cmake/
recipe-01-1.0.0-Linux/opt/recipe- 01/share/cmake/recipe-01/
recipe-01-1.0.0-Linux/opt/recipe- 01/share/cmake/recipe-01/messageConfig.cmake
recipe-01-1.0.0- Linux/opt/recipe-01/share/cmake/recipe-01/messageTargets-hello- world.cmake
recipe-01-1.0.0-Linux/opt/recipe-01/share/cmake/recipe- 01/messageConfigVersion.cmake
recipe-01-1.0.0-Linux/opt/recipe- 01/share/cmake/recipe-01/messageTargets-hello-world- release.cmake
recipe-01-1.0.0-Linux/opt/recipe-01/share/cmake/recipe- 01/messageTargets-release.cmake
recipe-01-1.0.0-Linux/opt/recipe- 01/share/cmake/recipe-01/messageTargets.cmake
recipe-01-1.0.0- Linux/opt/recipe-01/include/
recipe-01-1.0.0-Linux/opt/recipe- 01/include/message/
recipe-01-1.0.0-Linux/opt/recipe- 01/include/message/Message.hpp
recipe-01-1.0.0-Linux/opt/recipe- 01/include/message/messageExport.h
recipe-01-1.0.0-Linux/opt/recipe- 01/lib64/
recipe-01-1.0.0-Linux/opt/recipe- 01/lib64/libmessage.so
recipe-01-1.0.0-Linux/opt/recipe- 01/lib64/libmessage.so.1`
\end{codebox}

\mySmallsection{平台原生的二进制安装}

我们希望每个平台原生二进制安装程序的配置略有不同。可以在单个\texttt{CMakeCPack.cmake}中使用CPack管理这些差异，就像例子中做的那样。

对于GNU/Linux系统，配置了\texttt{DEB}和\texttt{RPM}生成器:

\begin{codebox}[cmake]
if(UNIX)
  if(CMAKE_SYSTEM_NAME MATCHES Linux)
    list(APPEND CPACK_GENERATOR "DEB")
    set(CPACK_DEBIAN_PACKAGE_MAINTAINER "robertodr")
    set(CPACK_DEBIAN_PACKAGE_SECTION "devel")
    set(CPACK_DEBIAN_PACKAGE_DEPENDS "uuid-dev")
    
    list(APPEND CPACK_GENERATOR "RPM")
    set(CPACK_RPM_PACKAGE_RELEASE "1")
    set(CPACK_RPM_PACKAGE_LICENSE "MIT")
    set(CPACK_RPM_PACKAGE_REQUIRES "uuid-devel")
  endif()
endif()
\end{codebox}

我们的示例依赖于UUID库，\texttt{CPACK\_DEBIAN\_PACKAGE\_DEPENDS}和\texttt{cpack\_rpm\_package\_require}选项允许指定，包和数据库中对其他包的依赖关系。可以使用dpkg和rpm程序分别分析生成的\texttt{.deb}和\texttt{.rpm}包的内容。

注意，\texttt{CPACK\_PACKAGING\_INSTALL\_PREFIX}也会影响这些包生成器：我们的包将安装到\texttt{/opt/recipe-01}。

CMake真正提供了跨平台和可移植构建系统的支持。下面将使用Nullsoft脚本安装系统(NSIS)创建一个安装程序:

\begin{codebox}[cmake]
if(WIN32 OR MINGW)
  list(APPEND CPACK_GENERATOR "NSIS")
  set(CPACK_NSIS_PACKAGE_NAME "message")
  set(CPACK_NSIS_CONTACT "robertdr")
  set(CPACK_NSIS_ENABLE_UNINSTALL_BEFORE_INSTALL ON)
endif()
\end{codebox}

如果在macOS上构建项目，将启用\texttt{Bundle packager}:

\begin{codebox}[cmake]
if(APPLE)
  list(APPEND CPACK_GENERATOR "Bundle")
  set(CPACK_BUNDLE_NAME "message")
  configure_file(${PROJECT_SOURCE_DIR}/cmake/Info.plist.in Info.plist @ONLY)
  set(CPACK_BUNDLE_PLIST ${CMAKE_CURRENT_BINARY_DIR}/Info.plist)
  set(CPACK_BUNDLE_ICON ${PROJECT_SOURCE_DIR}/cmake/coffee.icns)
endif()
\end{codebox}

macOS的示例中，需要为包配置属性列表文件，这是通过\texttt{configure\_file}实现的。\texttt{Info.plist}的位置和包的图标，这些都可以通过CPack的变量进行设置。

\begin{myInfoBox}
\textbf{NOTE}：可以在这里阅读，关于属性列表格式的更多信息:\texthref{https://en.wikipedia.org/wiki/Property\_list}
\end{myInfoBox}

\mySmallsection{更多信息}

对\texttt{ CMakeCPack.cmake }进行设置，要比列出CPack的配置选项简单的多，我们可以将\texttt{CPACK\_*}变量的每个生成器设置放在单独的文件中，比如\texttt{CMakeCPackOptions.cmake}，并将这些设置包含到\texttt{CMakeCPack.cmake}使用\texttt{set(CPACK\_PROJECT\_CONFIG\_FILE "\$\{PROJECT\_SOUsRCE\_DIR\}/CMakeCPackOptions.cmake")}将设置包含入\texttt{ CMakeCPack.cmake}中。还可以在CMake时配置该文件，然后在CPack时包含该文件，这为配置多格式包生成器提供了一种简洁的方法(参见\texthref{https://cmake.org/cmake/help/v3.6/module/CPack.html} )。

与CMake中的所有工具一样，CPack功能强大、功能多样，并且提供了更多的灵活性和选项。感兴趣的读者应该看官方文档的命令行界面CPack (https://cmake.org/cmake/help/v3.6/manual/cpack.1.html )手册页，如何使用CPack生成器打包相关项目的更多细节(https://cmake.org/cmake/help/v3.6/module/CPack.html )。


% --- recipe 2: 11.2-chinese.md ---
\section{11.2 通过PyPI发布使用CMake/pybind11构建的C++/Python项目}

\begin{myInfoBox}
\textbf{NOTE}：此示例代码可以在 \texthref{https://github.com/dev-cafe/cmake-cookbook/tree/v1.0/chapter-11/recipe-02} 中找到。该示例在CMake 3.11版(或更高版本)中是有效的，并且已经在GNU/Linux、macOS和Windows上进行过测试。
\end{myInfoBox}

本示例中，我们将以第9章第5节的代码的pybind11为例，为其添加相关的安装目标和pip打包信息，并将项目上传到PyPI。我们要实现一个可以使用pip安装，并运行CMake从而获取底层pybind11依赖项的项目。

\mySmallsection{准备工作}

要通过PyPI分发包的话，需要一个\texthref{https://pypi.org} 帐户。当然，也可以先从本地路径进行安装练习。

\begin{myInfoBox}
\textbf{TIPS}：建议使用Pipenv (https://docs.pipenv.org )或虚拟环境(https://virtualenv.pypa )安装这个包和其他的Python包。
\end{myInfoBox}

我们基于第9章第5节的项目，它包含一个主\texttt{CMakeLists.txt}文件和一个\texttt{account/CMakeLists.txt}文件，配置帐户示例目标时，使用如下的项目树：

\begin{codebox}[shell]
.
├── account
│    ├── account.cpp
│    ├── account.hpp
│    ├── CMakeLists.txt
│    s└── test.py
└── CMakeLists.txt
\end{codebox}

示例中，\texttt{account.cpp},\texttt{ account.hpp}和\texttt{test.py}没有任何变化。修改\texttt{account/CMakeLists.txt}，并为pip添加几个文件，以便能够构建安装包。为此，需要根目录中的另外三个文件：\texttt{README.rst}，\texttt{MANIFEST.in}和\texttt{setup.py}。

\texttt{README.rst}中包含关于项目的s文档：

\begin{codebox}[text]
Example project
===============

Project description in here ...
\end{codebox}

\texttt{MANIFEST.in}列出了需要安装的Python模块：

\begin{codebox}[text]
include README.rst CMakeLists.txt
recursive-include account *.cpp *.hpp CMakeLists.txt
\end{codebox}

最后，\texttt{setup.py}包含构建指令和安装项目的说明：

\begin{codebox}[python]
import distutils.command.build as _build
import os
import sys
from distutils import spawn
from distutils.sysconfig import get_python_lib
from setuptools import setup


def extend_build():
  class build(_build.build):
    def run(self):
      cwd = os.getcwd()
      if spawn.find_executable('cmake') is None:
        sys.stderr.write("CMake is required to build this package.\n")
        sys.exit(-1)
        _source_dir = os.path.split(__file__)[0]
        _build_dir = os.path.join(_source_dir, 'build_setup_py')
        _prefix = get_python_lib()
        try:
          cmake_configure_command = [
              'cmake',
              '-H{0}'.format(_source_dir),
              '-B{0}'.format(_build_dir),
              '-DCMAKE_INSTALL_PREFIX={0}'.format(_prefix),
          ]
          _generator = os.getenv('CMAKE_GENERATOR')
          if _generator is not None:
            cmake_configure_command.append('-
                                          G{0}'.format(_generator))
          spawn.spawn(cmake_configure_command)
          spawn.spawn(
                ['cmake', '--build', _build_dir, '--target', 'install'])
          os.chdir(cwd)
        except spawn.DistutilsExecError:
          sys.stderr.write("Error while building with CMake\n")
          sys.exit(-1)
          _build.build.run(self)
  return build

_here = os.path.abspath(os.path.dirname(__file__))

if sys.version_info[0] < 3:
  with open(os.path.join(_here, 'README.rst')) as f:
    long_description = f.read()
else:
  with open(os.path.join(_here, 'README.rst'), encoding='utf-8') as f:
    long_description = f.read()

_this_package = 'account'

version = {}
with open(os.path.join(_here, _this_package, 'version.py')) as f:
  exec(f.read(), version)
  
setup(
    name=_this_package,
    version=version['__version__'],
    description='Description in here.',
    long_description=long_description,
    author='Bruce Wayne',
    author_email='bruce.wayne@example.com',
    url='http://example.com',
    license='MIT',
    packages=[_this_package],
    include_package_data=True,
    classifiers=[
        'Development Status :: 3 - Alpha',
        'Intended Audience :: Science/Research',
        'Programming Language :: Python :: 2.7',
        'Programming Language :: Python :: 3.6'
    ],
    cmdclass={'build': extend_build()})
\end{codebox}

\texttt{account}子目录中放置一个\texttt{\_\_init\_\_.py}脚本：

\begin{codebox}[python]
from .version import __version__
from .account import Account
__all__ = [
  '__version__',
  'Account',
]
\end{codebox}

再放一个\texttt{version.py}脚本：

\begin{codebox}[python]
__version__ = '0.0.0'
\end{codebox}

项目的文件结构如下：

\begin{codebox}[shell]
.
├── account
│    ├── account.cpp
│    ├── account.hpp
│    ├── CMakeLists.txt
│    ├── __init__.py
│    ├── test.py
│    └── version.py
├── CMakeLists.txt
├── MANIFEST.in
├── README.rst
└── setup.py
\end{codebox}

\mySmallsection{具体实施}

本示例基于第9章第5节项目的基础上。

首先，修改\texttt{account/CMakeLists.txt}，添加安装目标：

\begin{codebox}[cmake]
install(
  TARGETS
    account
  LIBRARY
    DESTINATION account
  )
\end{codebox}

安装目标时，\texttt{README.rst}, \texttt{MANIFEST.in}，\texttt{setup.py}、\texttt{\_\_init\_\_.py}和\texttt{version.py}将放置在对应的位置上，我们准备使用pybind11测试安装过程：

\begin{enumerate}
\item 为此，在某处创建一个新目录，我们将在那里测试安装。

\item 在创建的目录中，从本地路径运行\texttt{pipenv install}。调整本地路径，指向\texttt{setup.py}的目录：

\begin{codebox}[shell]
$ pipenv install /path/to/cxx-example
\end{codebox}

\item 在Pipenv环境中打开一个Python shell：

\begin{codebox}[shell]
$ pipenv run python
\end{codebox}

\item Python shell中，可以测试我们的CMake包：

\begin{codebox}[shell]
>>> from account import Account
>>> account1 = Account()
>>> account1.deposit(100.0)
>>> account1.deposit(100.0)
>>> account1.withdraw(50.0)
>>> print(account1.get_balance())
150.0
\end{codebox}

\end{enumerate}

\mySmallsection{工作原理}

\texttt{\$\{CMAKE\_CURRENT\_BINARY\_DIR\}}目录包含编译后的\texttt{account.cpython-36m-x86\_64-linux-gnu.so}，这个动态库就是使用pybind11构建Python模块。但是请注意，它的名称取决于操作系统(本例中是64位Linux)和Python环境(本例中是Python 3.6)。\texttt{setup.py}s脚本将运行CMake，并根据所选的Python环境(系统Python，Pipenv或虚拟环境)将Python模块安装到正确的路径下。

不过，在安装模块时面临两个挑战：

\begin{itemize}
\item 名称可变

\item CMake外部设置路径

\end{itemize}

可以使用下面的安装目标来解决这个问题，将在setup.py中定义安装目标位置：

\begin{codebox}[cmake]
install(
  TARGETS
    account
  LIBRARY
    DESTINATION account
  )
\end{codebox}

指示CMake将编译好的Python模块文件安装到相对于安装目标位置的\texttt{account}子目录中(第10章中详细讨论了如何设置目标位置)。\texttt{setup.py}将通过设置\texttt{CMAKE\_INSTALL\_PREFIX}来设置安装位置，并根据Python环境指向正确的路径。

让我们看看\texttt{setup.py}如何实现的。自下而上来看一下脚本：

\begin{codebox}[python]
setup(
  name=_this_package,
  version=version['__version__'],
  description='Description in here.',
  long_description=long_description,
  author='Bruce Wayne',
  author_email='bruce.wayne@example.com',
  url='http://example.com',
  license='MIT',
  packages=[_this_package],
  include_package_data=True,
  classifiers=[
    'Development Status :: 3 - Alpha',
    'Intended Audience :: Science/Research',
    'Programming Language :: Python :: 2.7',
    'Programming Language :: Python :: 3.6'
  ],
  cmdclass={'build': extend_build()})
\end{codebox}

该脚本包含许多占位符，还包含一些自解释的语句。这里我们将重点介绍最后一个指令\texttt{cmdclass}。这个指令中，通过自定义\texttt{extend\_build}函数扩展默认的构建步骤。这个默认的构建步骤如下：

\begin{codebox}[python]
def extend_build():
  class build(_build.build):
    def run(self):
      cwd = os.getcwd()
      if spawn.find_executable('cmake') is None:
        sys.stderr.write("CMake is required to build this package.\n")
        sys.exit(-1)
        _source_dir = os.path.split(__file__)[0]
        _build_dir = os.path.join(_source_dir, 'build_setup_py')
        _prefix = get_python_lib()
        try:
          cmake_configure_command = [
              'cmake',
              '-H{0}'.format(_source_dir),
              '-B{0}'.format(_build_dir),
              '-DCMAKE_INSTALL_PREFIX={0}'.format(_prefix),
          ]
          _generator = os.getenv('CMAKE_GENERATOR')
          if _generator is not None:
            cmake_configure_command.append('-
                                          G{0}'.format(_generator))
          spawn.spawn(cmake_configure_command)
          spawn.spawn(
                ['cmake', '--build', _build_dir, '--target', 'install'])
          os.chdir(cwd)
        except spawn.DistutilsExecError:
          sys.stderr.write("Error while building with CMake\n")
          sys.exit(-1)
          _build.build.run(self)
  return build
\end{codebox}

首先，检查CMake是否可用。函数执行了两个CMake命令：

\begin{codebox}[python]
cmake_configure_command = [
    'cmake',
    '-H{0}'.format(_source_dir),
    '-B{0}'.format(_build_dir),
    '-DCMAKE_INSTALL_PREFIX={0}'.format(_prefix),
]
_generator = os.getenv('CMAKE_GENERATOR')
if _generator is not None:
  cmake_configure_command.append('-
                                G{0}'.format(_generator))
spawn.spawn(cmake_configure_command)
spawn.spawn(
      ['cmake', '--build', _build_dir, '--target', 'install'])
\end{codebox}

我们可以设置\texttt{CMAKE\_GENERATOR}环境变量来修改生成器。安装目录如下方式设置：

\begin{codebox}[shell]
_prefix = get_python_lib()
\end{codebox}

从安装目录的根目录下，通过\texttt{distutils.sysconfig}导入\texttt{get\_python\_lib}函数。\texttt{cmake --build \_build\_dir --target install}命令以一种可移植的方式，构建和安装我们的项目。使用\texttt{\_build\_dir}而不使用\texttt{build}的原因是，在测试本地安装时，项目可能已经包含了一个\texttt{build}目录，这将与新安装过程发生冲突。对于已经上传到PyPI的包，构建目录的名称并不会带来什么影响。

\mySmallsection{更多信息}

现在我们已经测试了本地安装，准备将包上传到PyPI。在此之前，请确保\texttt{setup.py}中的元数据(例如：项目名称、联系方式和许可协议信息)是合理的，并且项目名称没有与PyPI已存在项目重名。在上传到\texthref{https://pypi.org} 之前，先测试PyPI(https://test.pypi.org )上，进行上载和下载的尝试。

上传之前，我们需要在主目录中创建一个名为\texttt{.pypirc}的文件，其中包含(替换成自己的\texttt{yourusername}和\texttt{yourpassword })：

\begin{codebox}[text]
[distutils]account
index-servers=
  pypi
  pypitest

[pypi]
username = yourusername
password = yourpassword

[pypitest]
repository = https://test.pypi.org/legacy/
username = yourusername
password = yourpassword
\end{codebox}

我们将分两步进行。首先，我们在本地创建Release包：

\begin{codebox}[shell]
$ python setup.py sdist
\end{codebox}

第二步中，使用Twine上传生成的分布数据(我们将Twine安装到本地的Pipenv中):

\begin{codebox}[shell]
$ pipenv run twine upload dist/* -r pypitest

Uploading distributions to https://test.pypi.org/legacy/
Uploading yourpackage-0.0.0.tar.gz
\end{codebox}

下一步，从测试实例到，将包安装到一个隔离的环境中：

\begin{codebox}[python]
$ pipenv shell
$ pip install --index-url https://test.pypi.org/simple/ yourpackage
\end{codebox}

当一切正常，就将我们的包上传到了PyPI：

\begin{codebox}[shell]
$ pipenv run twine upload dist/* -r pypi
\end{codebox}


% --- recipe 3: 11.3-chinese.md ---
\section{11.3 通过PyPI发布使用CMake/CFFI构建C/Fortran/Python项目}

\begin{myInfoBox}
\textbf{NOTE}：此示例代码可以在 \texthref{https://github.com/dev-cafe/cmake-cookbook/tree/v1.0/chapter-11/recipe-03} 中找到，其中有一个C++和Fortran示例。该示例在CMake 3.5版(或更高版本)中是有效的，并且已经在GNU/Linux、macOS和Windows上进行过测试。
\end{myInfoBox}

基于第9章第6节的示例，我们将重用前一个示例中的构建块，不过这次使用Python CFFI来提供Python接口，而不是pybind11。这个示例中，我们通过PyPI共享一个Fortran项目，这个项目可以是C或C++项目，也可以是任何公开C接口的语言，非Fortran就可以。

\mySmallsection{准备工作}

项目将使用如下的目录结构：

\begin{codebox}[shell]
.
├── account
│    ├── account.h
│    ├── CMakeLists.txt
│    ├── implementation
│    │    └── fortran_implementation.f90
│    ├── __init__.py
│    ├── interface_file_names.cfg.in
│    ├── test.py
│    └── version.py
├── CMakeLists.txt
├── MANIFEST.in
├── README.rst
└── setup.py
\end{codebox}

主\texttt{CMakeLists.txt}文件和\texttt{account}下面的所有源文件(\texttt{account/CMakeLists.txt}除外)与第9章中的使用方式相同。\texttt{README.rst}文件与前面的示例相同。\texttt{setup.py}脚本比上一个示例多了一行(包含\texttt{install\_require =['cffi']}的那一行):

\begin{codebox}[python]
# ... up to this line the script is unchanged
setup(
  name=_this_package,
  version=version['__version__'],
  description='Description in here.',
  long_description=long_description,
  author='Bruce Wayne',
  author_email='bruce.wayne@example.com',
  url='http://example.com',
  license='MIT',
  packages=[_this_package],
  install_requires=['cffi'],
  include_package_data=True,
  classifiers=[
    'Development Status :: 3 - Alpha',
    'Intended Audience :: Science/Research',
    'Programming Language :: Python :: 2.7',
    'Programming Language :: Python :: 3.6'
  ],
  cmdclass={'build': extend_build()})
\end{codebox}

\texttt{MANIFEST.in}应该与Python模块和包一起安装，并包含以下内容:

\begin{codebox}[text]
include README.rst CMakeLists.txt
recursive-include account *.h *.f90 CMakeLists.txt
\end{codebox}

\texttt{account}子目录下，我们看到两个新文件。一个\texttt{version.py}文件，其为\texttt{setup.py}保存项目的版本信息：

\begin{codebox}[python]
__version__ = '0.0.0'
\end{codebox}

子目录还包含\texttt{interface\_file\_names.cfg.in}文件:

\begin{codebox}[cmake]
[configuration]
header_file_name = account.h
library_file_name = $<TARGET_FILE_NAME:account>
\end{codebox}

\mySmallsection{具体实施}

讨论一下实现打包的步骤：

\begin{enumerate}
\item 示例基于第9章第6节，使用Python CFFI扩展了\texttt{account/CMakeLists.txt}，增加以下指令:

\begin{codebox}[cmake]
file(
  GENERATE OUTPUT ${CMAKE_CURRENT_BINARY_DIR}/interface_file_names.cfg
  INPUT ${CMAKE_CURRENT_SOURCE_DIR}/interface_file_names.cfg.in
  )
set_target_properties(account
  PROPERTIES
    PUBLIC_HEADER "account.h;${CMAKE_CURRENT_BINARY_DIR}/account_export.h"
    RESOURCE "${CMAKE_CURRENT_BINARY_DIR}/interface_file_names.cfg"
  )
install(
  TARGETS
    account
  LIBRARY
    DESTINATION account/lib
  RUNTIME
    DESTINATION account/lib
  PUBLIC_HEADER
    DESTINATION account/include
  RESOURCE
    DESTINATION account
  )
\end{codebox}

安装目标和附加文件准备好之后，就可以测试安装了。为此，会在某处创建一个新目录，我们将在那里测试安装。

\item 新创建的目录中，我们从本地路径运行pipenv install。调整本地路径，指向\texttt{setup.py}脚本保存的目录:

\begin{codebox}[shell]
$ pipenv install /path/to/fortran-example
\end{codebox}

\item 现在在Pipenv环境中生成一个Python shell:

\begin{codebox}[shell]
$ pipenv run python
\end{codebox}

\item Python shell中，可以测试CMake包:

\begin{codebox}[shell]
>>> import account
>>> account1 = account.new()
>>> account.deposit(account1, 100.0)
>>> account.deposit(account1, 100.0)
>>> account.withdraw(account1, 50.0)
>>> print(account.get_balance(account1))
150.0
\end{codebox}

\end{enumerate}

\mySmallsection{工作原理}

使用Python CFFI和CMake安装混合语言项目的扩展与第9章第6节的例子相对比，和使用Python CFFI的Python包多了两个额外的步骤:

\begin{enumerate}
\item 需要\texttt{setup.py}s

\item 安装目标时，CFFI所需的头文件和动态库文件，需要安装在正确的路径中，具体路径取决于所选择的Python环境

\end{enumerate}

\texttt{setup.py}的结构与前面的示例几乎一致，唯一的修改是包含\texttt{install\_require =['cffi']}，以确保安装示例包时，也获取并安装了所需的Python CFFI。\texttt{setup.py}脚本会自动安装\texttt{\_\_init\_\_.py}和\texttt{version.py}。\texttt{MANIFEST.in }中的改变不仅有\texttt{README.rst}和CMake文件，还有头文件和Fortran源文件:

\begin{codebox}[text]
include README.rst CMakeLists.txt
recursive-include account *.h *.f90 CMakeLists.txt
\end{codebox}

这个示例中，使用Python CFFI和\texttt{setup.py}打包CMake项目时，我们会面临三个挑战:

\begin{itemize}
\item 需要将\texttt{account.h}和\texttt{account\_export.h}头文件，以及动态库复制到系统环境中Python模块的位置。

\item 需要告诉\texttt{\_\_init\_\_.py}，在哪里可以找到这些头文件和库。第9章第6节中，我们使用环境变量解决了这些问题，不过使用Python模块时，不可能每次去都设置这些变量。

\item Python方面，我们不知道动态库文件的确切名称(后缀)，因为这取决于操作系统。

\end{itemize}

让我们从最后一点开始说起：不知道确切的名称，但在CMake生成构建系统时是知道的，因此我们在\texttt{interface\_file\_names.cfg,in}中使用生成器表达式，对占位符进行展开：

\begin{codebox}[text]
[configuration]
header_file_name = account.h
library_file_name = $<TARGET_FILE_NAME:account>
\end{codebox}

输入文件用来生成\texttt{\$\{CMAKE\_CURRENT\_BINARY\_DIR\}/interface\_file\_names.cfg}：

\begin{codebox}[cmake]
file(
  GENERATE OUTPUT ${CMAKE_CURRENT_BINARY_DIR}/interface_file_names.cfg
  INPUT ${CMAKE_CURRENT_SOURCE_DIR}/interface_file_names.cfg.in
  )
\end{codebox}

然后，将两个头文件定义为\texttt{PUBLIC\_HEADER}(参见第10章)，配置文件定义为\texttt{RESOURCE}：

\begin{codebox}[cmake]
set_target_properties(account
  PROPERTIES
    PUBLIC_HEADER "account.h;${CMAKE_CURRENT_BINARY_DIR}/account_export.h"
    RESOURCE "${CMAKE_CURRENT_BINARY_DIR}/interface_file_names.cfg"
)
\end{codebox}

最后，将库、头文件和配置文件安装到\texttt{setup.py}定义的安装路径中:

\begin{codebox}[cmake]
install(
  TARGETS
    account
  LIBRARY
    DESTINATION account/lib
  RUNTIME
    DESTINATION account/lib
  PUBLIC_HEADER
    DESTINATION account/include
  RESOURCE
    DESTINATION account
  )
\end{codebox}

注意，我们为库和运行时都设置了指向\texttt{account/lib}的目标。这对于Windows很重要，因为动态库具有可执行入口点，因此我们必须同时指定这两个入口点。

Python包将能够找到这些文件，要使用\texttt{account/\_\_init\_\_.py }来完成：

\begin{codebox}[python]
# this interface requires the header file and library file
# and these can be either provided by interface_file_names.cfg
# in the same path as this file
# or if this is not found then using environment variables
_this_path = Path(os.path.dirname(os.path.realpath(__file__)))
_cfg_file = _this_path / 'interface_file_names.cfg'
if _cfg_file.exists():
  config = ConfigParser()
  config.read(_cfg_file)
  header_file_name = config.get('configuration', 'header_file_name')
  _header_file = _this_path / 'include' / header_file_name
  _header_file = str(_header_file)
  library_file_name = config.get('configuration', 'library_file_name')
  _library_file = _this_path / 'lib' / library_file_name
  _library_file = str(_library_file)
else:
  _header_file = os.getenv('ACCOUNT_HEADER_FILE')
  assert _header_file is not None
  _library_file = os.getenv('ACCOUNT_LIBRARY_FILE')
  assert _library_file is not None
\end{codebox}

本例中，将找到\texttt{\_cfg\_file}并进行解析，\texttt{setup.py}将找到\texttt{include}下的头文件和\texttt{lib}下的库，并将它们传递给CFFI，从而构造库对象。这也是为什么，使用\texttt{lib}作为安装目标\texttt{DESTINATION}，而不使用\texttt{CMAKE\_INSTALL\_LIBDIR}的原因(否则可能会让\texttt{account/\_\_init\_\_.py }混淆)。

\mySmallsection{更多信息}

将包放到PyPI测试和生产实例中的后续步骤，因为有些步骤是类似的，所以可以直接参考前面的示例。


% --- recipe 4: 11.4-chinese.md ---
\#11.4 以Conda包的形式发布一个简单的项目

\begin{myInfoBox}
\textbf{NOTE}：此示例代码可以在 \texthref{https://github.com/dev-cafe/cmake-cookbook/tree/v1.0/chapter-11/recipe-04} 中找到。该示例在CMake 3.5版(或更高版本)中是有效的，并且已经在GNU/Linux、macOS和Windows上进行过测试。
\end{myInfoBox}

虽然PyPI是发布Python包的标准平台，但Anaconda (https://anaconda.org )更为可能更为流行，因为它不仅允许使用Python接口发布Python或混合项目，还允许对非Python项目进行打包和依赖关系管理。这个示例中，我们将为一个非常简单的C++示例项目准备一个Conda包，该项目使用CMake配置和构建，除了C++之外没有依赖关系。下一个示例中，我们将会来看看一个更复杂的Conda包。

\mySmallsection{准备工作}

我们的目标是打包以下示例代码(\texttt{example.cpp})：

\begin{codebox}[cpp]
#include <iostream>
int main() {
    std::cout << "hello from your conda package!" << std::endl;
    return 0;
}
\end{codebox}

\mySmallsection{具体实施}

\begin{enumerate}
\item \texttt{CMakeLists.txt}文件给出了最低版本要求、项目名称和支持的语言:

\begin{codebox}[cmake]
cmake_minimum_required(VERSION 3.5 FATAL_ERROR)
project(recipe-04 LANGUAGES CXX)
set(CMAKE_CXX_STANDARD 11)
set(CMAKE_CXX_EXTENSIONS OFF)
set(CMAKE_CXX_STANDARD_REQUIRED ON)
\end{codebox}

\item 使用\texttt{example.cpp}构建\texttt{hello-conda}可执行目标：

\begin{codebox}[cmake]
add_executable(hello-conda "")
target_sources(hello-conda
  PRIVATE
    example.cpp
  )
\end{codebox}

\item 使用\texttt{CMakeLists.txt}定义安装目标：

\begin{codebox}[cmake]
nstall(
  TARGETS
    hello-conda
  DESTINATION
    bin
  )
\end{codebox}

\item 将在一个名为\texttt{ meta.yaml }的文件中，对Conda包进行描述。我们将把它放在\texttt{conda-recipe}目录下，文件结构如下：

\begin{codebox}[shell]
.
├── CMakeLists.txt
├── conda-recipe
│    └── meta.yaml
└── example.cpp
\end{codebox}

\item \texttt{meta.yaml}包含如下内容：

\begin{codebox}[yaml]
package:
  name: conda-example-simple
  version: "0.0.0"
source:
  path: .. /  # this can be changed to git-url
build:
  number: 0
  binary_relocation: true
  script:
    - cmake -H. -Bbuild_conda -G "${CMAKE_GENERATOR}" -DCMAKE_INSTALL_PREFIX=${PREFIX} # [not win]
    - cmake -H. -Bbuild_conda -G "%CMAKE_GENERATOR%" -DCMAKE_INSTALL_PREFIX="%LIBRARY_PREFIX%" # [win]
    - cmake - -build build_conda - -target install
requirements:
  build:
    - cmake >=3.5
    - { { compiler('cxx') } }
about:
  home: http://www.example.com
  license: MIT
  summary: "Summary in here ..."
\end{codebox}

\item 现在来构建包：

\begin{codebox}[shell]
$ conda build conda-recipe
\end{codebox}

\item 过程中屏幕上看到大量输出，但是一旦构建完成，就可以对包进行安装。首先，在本地进行测试：

\begin{codebox}[shell]
$ conda install --use-local conda-example-simple
\end{codebox}

\item 现在准备测试安装包，打开一个新的终端(假设Anaconda处于激活状态)，并输入以下内容：

\begin{codebox}[shell]
$ hello-conda
hello from your conda package!
\end{codebox}

\item 测试成功后，再移除包装：

\begin{codebox}[shell]
$ conda remove conda-example-simple
\end{codebox}

\end{enumerate}

\mySmallsection{工作原理}

\texttt{CMakeLists.txt}中，安装目标是这个示例的一个基本组件:

\begin{codebox}[cmake]
install(
  TARGETS
    hello-conda
  DESTINATION
    bin
  )
\end{codebox}

目标的二进制文件会安装到\texttt{\$\{CMAKE\_INSTALL\_PREFIX\}/bin}中。变量由Conda定义，并且构建步骤中定义在\texttt{ meta.yaml }：

\begin{codebox}[yaml]
build:
  number: 0
  binary_relocation: true
  script:
    - cmake -H. -Bbuild_conda -G "${CMAKE_GENERATOR}" -DCMAKE_INSTALL_PREFIX=${PREFIX} # [not win]
    - cmake -H. -Bbuild_conda -G "%CMAKE_GENERATOR%" -DCMAKE_INSTALL_PREFIX="%LIBRARY_PREFIX%" # [win]
    - cmake - -build build_conda - -target install
\end{codebox}

将安装目录设置为\texttt{\$\{prefix\}} (Conda的内部变量)，然后构建并安装项目。调用构建目录命名为\texttt{build\_conda}的动机与前面的示例类似：特定的构建目录名可能已经命名为\texttt{build}。

\mySmallsection{更多信息}

配置文件\texttt{ meta.yaml }可为任何项目指定构建、测试和安装步骤。详情请参考官方文档：\texthref{https://conda.io/docs/user-guide/tasks/build-packages/define-metadata.html}

要将Conda包上传到Anaconda云，请遵循官方的Anaconda文档: \texthref{https://docs.anaconda.com/anaconda-cloud/user-guide/}

此外，也可以考虑将Miniconda，作为Anaconda的轻量级替代品：\texthref{https://conda.io/miniconda.html}


% --- recipe 5: 11.5-chinese.md ---
\section{11.5 将Conda包作为依赖项发布给项目}

\begin{myInfoBox}
\textbf{NOTE}：此示例代码可以在 \texthref{https://github.com/dev-cafe/cmake-cookbook/tree/v1.0/chapter-11/recipe-05} 中找到。该示例在CMake 3.5版(或更高版本)中是有效的，并且已经在GNU/Linux、macOS和Windows上进行过测试。
\end{myInfoBox}

这个示例中，我们将基于之前示例的结果，并且为CMake项目准备一个更真实和复杂的Conda包，这将取决于DGEMM的函数实现，对于矩阵与矩阵的乘法，可以使用Intel的MKL库进行。Intel的MKL库可以以Conda包的形式提供。此示例将为我们提供一个工具集，用于准备和共享具有依赖关系的Conda包。

\mySmallsection{准备工作}

对于这个示例，我们将使用与前一个示例中的Conda配置，和相同的文件命名和目录结构：

\begin{codebox}[shell]
.
├── CMakeLists.txt
├── conda-recipe
│    └── meta.yaml
└── example.cpp
\end{codebox}

示例文件(\texttt{example.cpp})将执行矩阵-矩阵乘法，并将MKL库返回的结果与“noddy”实现进行比较:

\begin{codebox}[cpp]
#include "mkl.h"

#include <cassert>
#include <cmath>
#include <iostream>
#include <random>

int main() {
  // generate a uniform distribution of real number between -1.0 and 1.0
  std::random_device rd;
  std::mt19937 mt(rd());
  std:: uniform_real_distribution < double > dist(-1.0, 1.0);
  
  int m = 500;
  int k = 1000;
  int n = 2000;

  double *A = (double *)mkl_malloc(m * k * sizeof(double), 64);
  double *B = (double *)mkl_malloc(k * n * sizeof(double), 64);
  double *C = (double *)mkl_malloc(m * n * sizeof(double), 64);
  double * D = new double[m * n];
  
  for (int i = 0; i < (m * k); i++) {
    A[i] = dist(mt);
  }

  for (int i = 0; i < (k * n); i++) {
    B[i] = dist(mt);
  }

  for (int i = 0; i < (m * n); i++) {
    C[i] = 0.0;
  }

  double alpha = 1.0;
  double beta = 0.0;
  cblas_dgemm(CblasRowMajor,
              CblasNoTrans,
              CblasNoTrans,
              m,
              n,
              k,
              alpha,
              A,
              k,
              B,
              n,
              beta,
              C,
              n);

  // D_mn = A_mk B_kn
  for (int r = 0; r < m; r++) {
    for (int c = 0; c < n; c++) {
      D[r * n + c] = 0.0;
      for (int i = 0; i < k; i++) {
        D[r * n + c] += A[r * k + i] * B[i * n + c];
      }
    }
  }

  // compare the two matrices
  double r = 0.0;
  for (int i = 0; i < (m * n); i++) {
    r += std::pow(C[i] - D[i], 2.0);
  }
  assert (r < 1.0e-12 & & "ERROR: matrices C and D do not match");
  
  mkl_free(A);
  mkl_free(B);
  mkl_free(C);
  delete[] D;

  std:: cout << "MKL DGEMM example worked!" << std:: endl;
  
  return 0;`
}
\end{codebox}

我们还需要修改\texttt{meta.yaml}。然而，与上一个示例相比，唯一的变化是在依赖项中加入了\texttt{mkl-devel}：

\begin{codebox}[yaml]
package:
  name: conda-example-dgemm
  version: "0.0.0"

source:
  path: ../ # this can be changed to git-url

build:
  number: 0
  script:
  - cmake -H. -Bbuild_conda -G "${CMAKE_GENERATOR}"
  -DCMAKE_INSTALL_PREFIX=${PREFIX} # [not win]
  - cmake -H. -Bbuild_conda -G "%CMAKE_GENERATOR%"
  -DCMAKE_INSTALL_PREFIX="%LIBRARY_PREFIX%" # [win]
  - cmake - -build build_conda - -target install
  
requirements:
  build:
    - cmake >=3.5
    - {{ compiler('cxx') }}
  host:
    - mkl - devel 2018
    
about:
  home: http://www.example.com
  license: MIT
  summary: "Summary in here ..."
\end{codebox}

\mySmallsection{具体实施}

\begin{enumerate}
\item \texttt{CMakeLists.txt}文件声明了最低版本、项目名称和支持的语言：

\begin{codebox}[cmake]
cmake_minimum_required(VERSION 3.5 FATAL_ERROR)
project(recipe-05 LANGUAGES CXX)
set(CMAKE_CXX_STANDARD 11)
set(CMAKE_CXX_EXTENSIONS OFF)
set(CMAKE_CXX_STANDARD_REQUIRED ON)
\end{codebox}

\item 使用\texttt{example.cpp}构建\texttt{dgem-example}可执行目标：

\begin{codebox}[cmake]
add_executable(dgemm-example "")
target_sources(dgemm-example
  PRIVATE
    example.cpp
  )
\end{codebox}

\item 然后，需要找到通过\texttt{MKL-devel}安装的MKL库。我们准备了一个名为\texttt{IntelMKL}的\texttt{INTERFACE }库，该库可以用于其他目标，并将为依赖的目标设置包括目录、编译器选项和链接库。根据Intel的建议(https://software.intel.com/en-us/articles/intel-mml-link-line-advisor/ )进行设置。首先，设置编译器选项：

\begin{codebox}[cmake]
add_library(IntelMKL INTERFACE)
target_compile_options(IntelMKL
  INTERFACE
    $<$<OR:$<CXX_COMPILER_ID:GNU>,$<CXX_COMPILER_ID:AppleClang>>:-m64>
  )
\end{codebox}

\item 接下来，查找\texttt{mkl.h}头文件，并为\texttt{IntelMKL}目标设置\texttt{include}目录：

\begin{codebox}[cmake]
find_path(_mkl_h
  NAMES
    mkl.h
  HINTS
    ${CMAKE_INSTALL_PREFIX}/include
  )
target_include_directories(IntelMKL
  INTERFACE
    ${_mkl_h}
  )
message(STATUS "MKL header file FOUND: ${_mkl_h}")
\end{codebox}

\item 最后，为\texttt{IntelMKL}目标设置链接库:

\begin{codebox}[cmake]
find_library(_mkl_libs
  NAMES
    mkl_rt
  HINTS
    ${CMAKE_INSTALL_PREFIX}/lib
  )
message(STATUS "MKL single dynamic library FOUND: ${_mkl_libs}")
find_package(Threads QUIET)
target_link_libraries(IntelMKL
  INTERFACE
    ${_mkl_libs}
    $<$<OR:$<CXX_COMPILER_ID:GNU>,$<CXX_COMPILER_ID:AppleClang>>:Threads::Threads>
    $<$<OR:$<CXX_COMPILER_ID:GNU>,$<CXX_COMPILER_ID:AppleClang>>:m>
  )
\end{codebox}

\item 使用\texttt{cmake\_print\_properties}函数，打印\texttt{IntelMKL}目标的信息：

\begin{codebox}[cmake]
include(CMakePrintHelpers)
cmake_print_properties(
  TARGETS
    IntelMKL
  PROPERTIES
    INTERFACE_COMPILE_OPTIONS
    INTERFACE_INCLUDE_DIRECTORIES
    INTERFACE_LINK_LIBRARIES
  )
\end{codebox}

\item 将这些库连接到\texttt{dgem-example}:

\begin{codebox}[cmake]
target_link_libraries(dgemm-example
  PRIVATE
    IntelMKL
  )
\end{codebox}

\item \texttt{CMakeLists.txt}中定义了安装目标:

\begin{codebox}[cmake]
install(
  TARGETS
    dgemm-example
  DESTINATION
    bin
  )
\end{codebox}

\item 尝试构建包：

\begin{codebox}[shell]
$ conda build conda-recipe
\end{codebox}

\item 过程中屏幕上将看到大量输出，但是一旦构建完成，就可以对包进行安装包。首先，在本地进行安装测试：

\begin{codebox}[shell]
$ conda install --use-local conda-example-dgemm
\end{codebox}

\item 现在测试安装，打开一个新的终端(假设Anaconda处于激活状态)，并输入：

\begin{codebox}[shell]
$ dgemm-example
MKL DGEMM example worked!
\end{codebox}

\item 安装成功之后，再进行卸载：

\begin{codebox}[shell]
$ conda remove conda-example-dgemm
\end{codebox}

\end{enumerate}

\mySmallsection{工作原理}

\texttt{meta.yaml }中的变化就是\texttt{mml-devel}依赖项。从CMake的角度来看，这里的挑战是定位Anaconda安装的MKL库。幸运的是，我们知道它位于\texttt{\$\{CMAKE\_INSTALL\_PREFIX\}}中。可以使用在线的\texttt{Intel MKL link line advisor}(https://software.intel.com/en-us/articles/intel-mml-link-line-advisor/) 查看如何根据选择的平台和编译器，将MKL链接到我们的项目中，我们会将此信息封装到\texttt{INTERFACE}库中。这个解决方案非常适合类MKL的情况：库不是由我们的项目或任何子项目创建的目标，但是它仍然需要以一种方式进行处理；也就是：设置编译器标志，包括目录和链接库。\texttt{INTERFACE}库是构建系统中的目标，但不创建任何构建输出(至少不会直接创建)。但由于它们是目标，我们可对它们的属性进行设置。这样与“实际”目标一样，可以安装、导出和导入。

首先，我们用\texttt{INTERFACE}属性声明一个名为\texttt{IntelMKL}的新库。然后，根据需要设置属性，并使用\texttt{INTERFACE}属性在目标上调用适当的CMake命令：

\begin{itemize}
\item target\_compile\_options：用于设置\texttt{INTERFACE\_COMPILE\_OPTIONS}。示例中，设置了\texttt{-m64}，不过这个标志只有GNU和AppleClange编译器能够识别。并且，我们使用生成器表达式来实现。

\item target\_include\_directories：用于设置\texttt{INTERFACE\_INCLUDE\_DIRECTORIES}。使用\texttt{find\_path}，可以在找到系统上的\texttt{mkl.h}头文件后设置这些参数。

\item target\_link\_libraries：用于设置\texttt{INTERFACE\_LINK\_LIBRARIES}。我们决定链接动态库\texttt{ libmkl\_rt.so }，并用\texttt{find\_library}搜索它。GNU或AppleClang编译器还需要将可执行文件链接到线程和数学库。同样，这些情况可以使用生成器表达式优雅地进行处理。

\end{itemize}

在\texttt{IntelMKL}目标上设置的属性后，可以通过\texttt{cmake\_print\_properties}命令将属性进行打印。最后，链接到\texttt{IntelMKL}目标，这将设置编译器标志，包括目录和链接库：

\begin{codebox}[cmake]
target_link_libraries(dgemm-example
  PRIVATE
    IntelMKL
  )
\end{codebox}

\mySmallsection{更多信息}

Anaconda云上包含大量包。使用上述方法，可以为CMake项目构建依赖于其他Conda包的Conda包。这样，就可以探索软件功能的各种可能性，并与他人分享您的软件包!


